\section{Language}
\label{sec:formal}

\begin{figure}[th!]
  \vspace{-.3cm}
{\footnotesize
  \begin{alignat*}{2}
    \text{\textbf{Variables }}& \quad &\quad& x, y, z, \vnu, ...
    \\[-0.2em]\text{\textbf{Base Types}}& \quad & b  ::= \quad &  \Code{unit} ~|~ \Code{bool} ~|~ \Code{nat} ~|~ \Code{int} ~|~ \Code{id} ...
    \\[-0.2em]\text{\textbf{Basic Types}}& \quad & s  ::= \quad &  b ~|~ s\sarr s
    \\\text{\textbf{Pure Operations}}& \quad & \primop ::= \quad & {+} ~|~ {-} ~|~ {==} ~|~ {<} ~|~ {\leq} ~|~ \Code{rand\_int}~...
    \\[-0.2em]\text{\textbf{Constants }}& \quad &c ::= \quad & () ~|~ \mathbb{B} ~|~ \mathbb{Z} ~|~ \mathbb{I}...
    \\[-0.2em]\text{\textbf{Qualifiers}}& \quad & \phi ::= \quad & c ~|~ x ~|~ \primop\;\overline{v} ~|~ \bot ~|~ \top  ~|~ \neg \phi ~|~ \phi \land \phi ~|~ \phi \lor \phi ~|~ \phi \impl \phi ~|~ \forall x{:}b.\, \phi
    \\[-0.2em]\text{\textbf{Effectful Operations}}& \quad & \eff{op} ::= \quad & \eff{put} ~|~ \eff{get} ~|~ \eff{readReq} ~|~ \eff{readRsp} ~|~ \eff{write} ~|~ \eff{push} ~|~ \eff{pop} ~|~ ...
    % \\\text{\textbf{Event Kinds}}& \quad & k ::= \quad & \eff{gen}
    % ~|~ \eff{obs}
      % ~|~ \reqOp{op}{\ensuremath{\overline{v}}} ~|~ \respOp{op}{\ensuremath{\iota}}
    \\[-0.2em]\text{\textbf{Values }}& \quad & v ::= \quad & c ~|~ x
    ~|~ \zlam{x}{s}{e} ~|~ \zfix{f}{s}{x}{s}{e}
    \\[-0.2em]\text{\textbf{Expressions}}& \quad & e ::=\quad & v ~|~
    e \oplus e ~|~ v\; v ~|~ \primop\; \overline{v} ~|~ \eff{op}\;
    \overline{v} ~|~\zlet{x{:}b}{e}{e} ~|~ \sassume{\phi} ~|~
    \sassert{\phi}
    \\[-0.2em]\text{\textbf{Events}} & \quad & m ::= \quad &
    \langle\eff{op}{\; \overline{v}}\rangle \\
    \text{\textbf{Traces}} & \quad & \alpha ::= \quad & \emptr ~|~ m
    \;{::}\; \alpha ~|~ \alpha \listconcat \alpha
    \\[5pt]
    \text{\textbf{Notations}} & \quad & e_1;e_2 \doteq \quad &
    \zlet{\_}{e_1}{e_2}  \\
    & & \hspace{-2.5cm}\zlet{x{:}t}{\sassume{\phi(x)}}{e} \doteq \quad &
    \zlet{x{:}t}{\Code{rand\_int ()}}{\sassume{\phi(x)}; e}
    %  \mykw{let}\;\eff{op}\;\overline{v}\;\mykw{in}\;e \doteq \zgen{op}{\overline{v}}{e} \qquad\qquad
  \end{alignat*}
}
\vspace*{-.3in}
\caption{Syntax of \DSL{} expressions}
\label{fig:term-syntax}
\vspace*{-.2in}
\end{figure}

We formalize our approach using a core language, \DSL{}, for
expressing \textit{trace generators}-- nondeterministic programs that
probe the behaviors of a black-box system under test via a set of
(effectful) operators that the system provides. This calculus is
equipped with a type system that characterizes the set of traces that
a program \emph{must} be able to produce, \emph{if} the system under
test produces appropriate responses. The syntax of \DSL{} is shown in
\autoref{fig:term-syntax}; the calculus includes both pure and
effectful operations ($\primop$ and $\eff{op}$), nondeterministic
choices ($\oplus$), recursion ($\mykw{fix}$), and $\mykw{assume}$ and
$\mykw{assert}$ statements in the style of solver-aided
languages~\cite{BC+10}. As we saw in \autoref{sec:overview},
asynchronous effects can be modeled as a pair of $\eff{opReq}$ /
$\eff{opResp}$ operations, where the former generates an id (given as
$\iota$ in the syntax) that can then be subsequently referenced when
performing a corresponding response operation.
%\autoref{fig:term-syntax} defines some shorthand notations used by our
%examples, e.g., we use $e_1;e_2$ to simplify $\mykw{let}$-expressions
%when the remaining program does not use the bound variable, i.e.,
%$\zlet{\_}{e_1}{e_2}$.

\begin{wrapfigure}{r}{.5\textwidth}
  \footnotesize
  \begin{minipage}{.45\textwidth}
    \begin{flalign*}
    &\text{\textbf{Operator Semantics }}\
    \fbox{$\vDash \primop(\overline{c}) \Downarrow c \quad \alpha \vDash \eff{op}(\overline{c}) \Downarrow c$} \\
    & \text{\textbf{Operational Semantics }}\
    \fbox{$
    % \phi \Downarrow c \quad
    \steptr{\alpha}{e}{\alpha}{e}$}
\end{flalign*}
\end{minipage}
\\[5pt]
\begin{center}
  \begin{prooftree}
    \hypo{\vDash \primop(\overline{c}) \Downarrow c'}
    \infer1[\textsc{\small StPOp}]{
      \steptr{\alpha}{\primop \; \overline{c}\;}{[\;]}{c'}
    }
  \end{prooftree}
  \begin{prooftree}
    \hypo{\alpha \vDash \eff{op}\overline{c} \Downarrow c'}
    \infer1[\textsc{\small StEfOp}]{
      \steptr{\alpha}{\eff{op}\; \overline{c}\;}{[\eff{op}\;\overline{c}\;c']}{c'}
    }
  \end{prooftree}
  \end{center}
  \vspace*{-.05in}
    \caption{Selected reduction rules}
    \label{fig:selected-semantics}
  \vspace*{-.2in}
\end{wrapfigure}
\DSL{} is equipped with a small-step operational semantics similar to
other calculi for trace-based type systems~\cite{ZYDJ24}; selected
reduction rules are shown in
\autoref{fig:selected-semantics}.\footnote{The full set of reduction
  rules are included in the \techreport{}.} The reduction relation
$\steptr{\alpha}{e}{~\alpha'}{e'}$ constrains the output trace
$\alpha'$ produced by executing $e$ based on a \textit{history trace}
$\alpha$.  The evaluation of pure operations (rule \textsc{\small
  StPureOp}) make no use of trace, either history or output. The
result of an effectful operation, in contrast, can depend on the
events that preceded it, and it produces an output trace containing an
event whose payload contains both the arguments and return value of
the operation.

% \begin{example}[Auxiliary Operator Semantics]\label{ex:aux-semantics}
% The auxiliary semantics of the $\eff{put}$ and $\eff{get}$ operators in Example~\ref{ex:sc} can be defined as follows:

% {\footnotesize
%   \mprooftr{42mm}{}{OpPut}{
%   $\alpha \vDash \zevent{putReq}{k,v} \Downarrow \{ \zevent{putRsp}{k,v} \}$
%   }
%   \quad
%   \mprooftr{62mm}{}{OpGet}{
%   $\alpha_1\listconcat\zevent{putRsp}{k,v}\listconcat\alpha_2\vDash \zevent{getReq}{k} \Downarrow \{\zevent{getRsp}{k,v}\}$
%   }
%   \\[-0.1em]
%   }\noindent

% \noindent
% The database handles the $\eff{putReq}$ event by simply adding a $\eff{putRsp}$ event with the same payload to the capability. To guarantee the EC, handling a $\eff{getReq}$ adds a $\eff{getRsp}$ event where the read value is the one previously written value $v$ at the time the $\eff{getReq}$ occurs.
% \end{example}

% \begin{example}[Operational Semantics]\label{ex:semantics-stlc}
% We show a program that generates STLC term $\lambda.[0]$.
%   {\footnotesize
%     \begin{flalign*}
%       &[] \vDash
%       (\emptyset, \zassume{x}{\top}{\eff{abs}\;x;\eff{var}\;0;\mykw{obs}\;\eff{endAbs}})
%       \xhookrightarrow{[]} (\emptyset, \eff{abs}\;\Int;\eff{var}\;0;\mykw{obs}\;\eff{endAbs}) \tag{\textsc{StAssume}} &&
%       \\[-0.4em]
%       &\;\;\xhookrightarrow{[\zevent{abs}{\Int}]} (\{\eff{endAbs}\}, \eff{var}\;0;\mykw{obs}\;\eff{endAbs})  \xhookrightarrow{[\zevent{var}{0}]} (\{\eff{endAbs}\}, \mykw{obs}\;\eff{endAbs})
%       \xhookrightarrow{[\eff{endAbs}]} (\emptyset, ()) \tag{\textsc{StGen, StObs}} &&
%     \end{flalign*}
%   }\noindent
%     The program first instantiates $x$ with an arbitrary STLC type (e.g., $\Int$), applies $\eff{abs}$ to produce capability $\eff{endAbs}$, executes $\eff{var}$ which generates no new capability, and finally consumes $\eff{endAbs}$.
% \end{example}

% \begin{example}[Strong consistency violation]\label{ex:semantics-async}
%   The following program shows a strong consistency violation in Example~\ref{ex:sc}, under the auxiliary semantics of $\eff{put}$ and $\eff{get}$ shown in Example~\ref{ex:aux-semantics}.
%     {\footnotesize
%       \begin{flalign*}
%         &[\zevent{putReq}{``\Code{k}", 1};\zevent{putRsp}{``\Code{k}", 1}] \vDash
%         (\emptyset, \eff{putReq}\;``\Code{k}"\;2;\eff{getReq}\;``\Code{k}";\sobs{putRsp};\zobs{k, v}{getRsp}{\sassert{v \not= 2}})
%         \\[-0.2em]
%         &\;\;\xhookrightarrow{[\zevent{putReq}{``\Code{k}", 2}]} (\{ \zevent{putRsp}{``\Code{k}", 2} \}, \eff{getReq}\;``\Code{k}";\sobs{putRsp};\zobs{k, v}{getRsp}{\sassert{v \not= 2}})  \tag{\textsc{StGen}} &&
%         \\[-0.4em]
%         &\;\;\xhookrightarrow{[\zevent{getReq}{``\Code{k}"}]} (\{ \zevent{putRsp}{``\Code{k}", 2}, \zevent{getRsp}{``\Code{k}", 1} \}, \sobs{putRsp};\zobs{k, v}{getRsp}{\sassert{v \not= 2}})  \tag{\textsc{StGen}} &&
%         \\[-0.4em]
%         &\;\;\xhookrightarrow{[\zevent{putRsp}{``\Code{k}", 2}]} (\{\zevent{getRsp}{``\Code{k}", 1}\}, \zobs{k, v}{getRsp}{\sassert{v \not= 2}})  \tag{\textsc{StObs}} &&
%         \\[-0.4em]
%         &\;\;\xhookrightarrow{[\zevent{getRsp}{``\Code{k}", 1}]} (\emptyset, \sassert{1 \not= 2})  \tag{\textsc{StObs}} &&
%       \end{flalign*}
%     }\noindent
%   As shown above, the program update the value of key $\Code{k}$ to $2$, but the read value is still $1$ due to the auxiliary semantics of $\eff{getReq}$ with EC is allowed to return any previous written value.
% \end{example}

% We show a program that generates the identity function $\lambda.[0]$ in the style of \autoref{fig:spec} as follows:
% {\footnotesize
%   \begin{flalign*}
%     &[] \vDash
%     (\emptyset, \zassume{x}{\top}{\eff{abs}\;x;\eff{var}\;0;\eff{putTy}\;x;\eff{ayncGetTy};\zlet{y}{\mykw{obs}\;\eff{getTy}}{\mykw{obs}\;\eff{endAbs};\eff{putTy}\;\Code{Arr}(x, y)}})\\[-0.4em]
%     &\;\;\xhookrightarrow{[]} (\emptyset, \eff{abs}\;\Int;\eff{var}\;0;\eff{putTy}\;\Int;\eff{ayncGetTy};\zlet{y}{\mykw{obs}\;\eff{getTy}}{\mykw{obs}\;\eff{endAbs};\eff{putTy}\;\Code{Arr}(\Int, y)}) \tag{\textsc{StAssume}} &&
%     \\[-0.4em]
%     &\;\;\xhookrightarrow{[\zevent{abs}{\Int}]} (\{\eff{endAbs}\}, \eff{var}\;0;\eff{putTy}\;\Int;\eff{ayncGetTy};\zlet{y}{\mykw{obs}\;\eff{getTy}}{\mykw{obs}\;\eff{endAbs};\eff{putTy}\;\Code{Arr}(\Int, y)})  \tag{\textsc{StGen}} &&
%     \\[-0.6em]
%     &\;\;\xhookrightarrow{[\zevent{var}{0};\zevent{putTy}{\Int};\eff{ayncGetTy}]}^* (\{\eff{endAbs};\zevent{getTy}{\Int}\}, \zlet{y}{\mykw{obs}\;\eff{getTy}}{\mykw{obs}\;\eff{endAbs};\eff{putTy}\;\Code{Arr}(\Int, y)}) \tag{\textsc{StGen}} &&
%     \\[-0.6em]
%     &\;\;\xhookrightarrow{[\zevent{getTy}{\Int}]} (\{\eff{endAbs}\}, \mykw{obs}\;\eff{endAbs};\eff{putTy}\;\Code{Arr}(\Int, \Int)) \tag{\textsc{StObs}} &&
%     \\[-0.4em]
%     &\;\;\xhookrightarrow{[\eff{endAbs}]} (\emptyset, \eff{putTy}\;\Code{Arr}(\Int, \Int)) \xhookrightarrow{[\zevent{putTy}{\Code{Arr}(\Int, \Int)}]} (\emptyset, ()) \tag{\textsc{StObs, StGen}} &&
%   \end{flalign*}
% }\noindent
%   The first step substitutes $x$ with an arbitrary STLC type (e.g., $\Int$). In the next steps, the program performs the generative operator $\eff{abs}$, which generates a new capability $\eff{endAbs}$ to be handled in the future. Similarly, the program performs $\eff{var}$, $\eff{putTy}$, and $\eff{ayncGetTy}$, where only the last operator generates a new capability, $\zevent{getTy}{\Int}$, whose payload is the last $\eff{put}$ type.
%   The fourth step consumes the $\zevent{getTy}{\Int}$ event and substitutes $y$ with $\Int$ in the remaining program. The rest of the program is then executed, which ends the lambda abstraction and records the type of the generated STLC term as $\Code{Arr}(\Int, \Int)$.

\subsection{Type System}

The syntax of types in \DSL{} is shown in \autoref{fig:type-syntax}.
Types include \emph{pure} refinement types, which describe pure
computations, and underapproximate Hoare Automata Types (\tyName{}s),
which describe effectful computations.
% which describe incorrectness triple $\effLR{H}\;e\;\effLR{F}$ as a
% type $\uhat{H}{\Unit}{F}$ of term $e$.
Pure refinement types are similar to those found in prior
underapproximate refinement type systems~\cite{CoverageType}, and
allow base types (e.g., {\Code{int}}) to be further constrained by a
logical formula or qualifier. In contrast to traditional type systems
where a pure type qualifier $\ort{b}{\phi}$ constrains every value a
pure expresion \emph{may} evaluate to; the type qualifier
$\urt{b}{\phi}$ constrains a subset of the values it \emph{must}
evaluate to.  For example, the constant value $1$ can be assigned the
types $\ort{\Int}{\vnu = 1}$ and $\ort{\Int}{0 < \vnu}$ -- the value 1
satisfies the constraint in both \emph{overapproximate} qualifiers --
but, 1 cannot be assigned the type,
$\urt{\Int}{\vnu = 1 \lor \vnu = 2}$, since it cannot evaluate to both
1 and 2 as required by this \emph{underapproximate} qualifier. The
term $1 \oplus 2 \oplus 3$, in contrast, does have this
underapproximate type.

\paragraph{Underapproximate Hoare Automata Types}

Similar to other recent trace-based type systems~\cite{ZYDJ24}, \DSL{}
uses \emph{Symbolic Finite Automata} (SFAs)~\cite{Vea13,
  SFA-minimization, SFA-Transducers} to qualify the traces generated
by effectful computations; in contrast to those prior works, the
\tyName{}s in $\DSL{}$ use SFAs to \emph{underapproximate} the set of effects
that a program must produce. A \tyName{} $\uhat{H}{x{:}t}{F}$ consists
of a pure refinement type $t$ and two SFAs: a \emph{history} SFA $H$
that constrains the history trace under which a term is executed, and a
\emph{future} SFA $A$ that describes a subset of the events that it
must generate as a result. \tyName{}s represents SFAs using symbolic
regular expressions (SREs), but in principle could support any
encoding of SFAs, e.g., Symbolic LTL$_f$~\cite{LTLf}. SREs support
\emph{symbolic events}, $\msgB{op}{\overline{x}}{\phi}$, atomic
predicates that describes an effectful operation {$\eff{op}$} whose
arguments $\overline{x}$ satisfy the qualifier $\phi$. SREs provide a
convenient language for writing SFAs, supporting versions of the
standard regex operators shown in \autoref{fig:type-syntax}, including
union $A \lorA A$, sequence $A \seqA A$, and Kleene star $A^*$
operators.

\begin{figure}[tb!]
{\small
\begin{alignat*}{2}
    % \\\text{\textbf{Symbolic Event}}& \quad &se  ::= \quad & \msgB{op}{\overline{x}}{\phi}
    \text{\textbf{Symbolic Regex}}& \quad &H,F,A  ::= \quad & \emptyset ~|~ \epsilon ~|~ \msgB{op}{\overline{x}}{\phi} ~|~ A \lorA A ~|~ A \seqA A ~|~ A^* ~|~ A \setminus A ~|~ A \landA A ~|~ \anyA
    % \\& \quad &  \quad & \text{\textcolor{gray}{in Symbolic LTL$_f$ }}  se ~|~ A \lorA A ~|~ \nextA A ~|~ A \untilA A ~|~ \negA A ~|~ A \landA A
    \\\text{\textbf{Pure Refinement Types}}& \quad & t  ::= \quad & \urt{b}{\phi} ~|~  x{:}t\sarr t ~|~ x{:}t\sarr \tau ~|~ x{:}b \garr t ~|~ \effGray{op}{:}s \sarr t
    \\ \text{\textbf{Underapproximate HATs}}& \quad & \tau  ::= \quad & \uhat{H}{x{:}t}{F} ~|~ \tau \interty \tau
    \\\text{\textbf{Type Contexts}}& \quad &\Gamma ::= \quad  &\emptyset ~|~ x{:}t, \Gamma
  \end{alignat*}
  \vspace*{-.15in}
  {\footnotesize\begin{flalign*}
    {\small\text{\textbf{Type Aliases}}} & \;\; \msg{op}{\phi} \doteq \msgB{op}{\overline{x}}{\phi} \;\; \msgA{op}{\overline{v}} \doteq \msgB{op}{\overline{x}}{\overline{x = v}} \;\; \msgC{op} \doteq \msgB{op}{\overline{x}}{\top} \;\; A_1 A_2 \doteq A_1 \seqA A_2 \;\; b \doteq \urt{b}{\top}
    % \;\; \tau \doteq x{:}b\garr\tau
    % \\
    % & \widehat{\ort{b}{\phi}} \doteq \urt{b}{\phi} \;\; \widehat{\urt{b}{\phi}} \doteq \ort{b}{\phi} \;\; \widehat{x{:}t\sarr t'} \doteq x{:}t\sarr t' \;\; \widehat{x{:}t\sarr \tau} \doteq x{:}t\sarr \tau
  \end{flalign*}}
}
\vspace*{-.2in}
\caption{\DSL{} types}
\label{fig:type-syntax}
\vspace*{-.2in}
\end{figure}

The figure also gives the definition of the type aliases that
\autoref{sec:overview} used to limit user annotation burden.  When the
fields of a symbolic event are clear from context, we omit its
arguments $\overline{x}$, e.g., $\langle\eff{push}\;|\;v > 0\rangle$
means $\langle\eff{push}\;v\;|\;v > 0\rangle$. We omit quantifiers
with equality constraints, e.g., writing
$\langle\eff{push}\;v\;|\;v = 3\rangle$, as $\msgA{push}{3}$, and the
top ($\top$) qualifier in symbolic events and types.



%% \BD{We need to describe ghost variables at some point. If we are not
%%   going to use them in the intro+overview, we should explain that we
%%   left them are implicit there.}

\begin{example}[Ghost Events]
  \label{ex:stack-example}
  Consider an effectful stack library implemented as a fixed-size
  array whose elements grow monotonically.  An incorrect implementation
  that fails to grow the array when the stack is full may cause
  $\eff{push}$ operations to be dropped.  To detect this violation, we
  need to explore traces of the following form, where the last pop
  operation raises an ``empty stack'' exception:
  \vspace{-.6cm} \par\nobreak {\small%
    \begin{align*}
      \zevent{push}{1};\zevent{push}{2}; \zevent{push}{3}; \zevent{pop}{2}; \zevent{pop}{1};\textcolor{red}{\zevent{pop}{?}}
    \end{align*}}\noindent
  Because we would like the \tyName{} specifications for $\eff{push}$
  and $\eff{pop}$ to guide the synthesis of a test generator, they
  should be interpreted as underapproximations, consisting of nested
  pairs of $\eff{push}$ and $\eff{pop}$ events. This can be expressed
  by the following \tyName{} specifications, using ghost events
  $\effGray{pushI}$ and $\effGray{popI}$ as shown below:
  \vspace{-.4cm} \par\nobreak {\small %
    \begin{align*} %
      &\eff{push} :~ n{:}\Int \garr y{:}\Int \garr
        x{:}\urt{\Int}{y < \vnu} \sarr
        \uhat{\Code{Last}(\msgAGray{pushI}{n\;y})}{\Unit}{\msgA{push}{x}\seqA\msgAGray{pushI}{n{+}1\;x}}
      \\
      &\eff{pop} :~ n{:}\Int \garr m{:}\Int \garr \\[-0.4em]
      & \quad
        \uhat{\Code{Last}(\msgAGray{popI}{m\;\_}) \land
        \Code{Last}(\msgAGray{pushI}{n\;\_}) \land
        (\allA\seqA\msgAGray{pushI}{(n{-}m)\;x}\seqA\allA)}{x{:}\Int}{\msgA{pop}{x}\seqA\msgAGray{popI}{m{+}1\;
        x}}
    \end{align*}}\vspace{-.5cm} \par\nobreak\noindent
  Intuitively, the $\effGray{pushI}$ and $\effGray{popI}$ ghost events are
  simply $\eff{push}$ and $\eff{pop}$ events equipped with indices that
  record the number of push and pop operations.  As seen here, ghost events allow \tyName{}
    specifications to capture latent data dependencies between effectful
    operations.   The signature for $\eff{push}$ requires
  that the new pushed value ($x$) is greater than the last pushed value
  ($y$), and adds a new ghost event $\effGray{pushI}$ with an
  incremented index (i.e., $\msgAGray{pushI}{n{+}1\;x}$).
  The
  signature for $\eff{pop}$ also records the number of pop operations
  in the same way (i.e., $\msgAGray{popI}{m{+}1\;x}$); additionally, it
  requires the last popped value ($x$) to be equal to the
  $(n{-}m)^{th}$ pushed value, reflecting the ``last-in-first-out’'
  property of the stack.
\end{example}
%%\end{example}
%% \tyName{}s can only appear in the result type of a function, where
%% they describe the effects performed when the function is evaluated.
%% Following prior work~\cite{ZYDJ24}, our type system includes
%% \emph{ghost variables} ($x{:}b \garr \tau$), allowing it to explicitly
%% capture data dependencies between symbolic events in \tyName{}s. It
%% also includes analogous \emph{ghost events} ($\effGray{op}{:}s \sarr
%% \tau$, which are written in $\effGray{gray}$ to distinguish them from
%% regular effect operators), enabling \tyName{}s to describe non-regular
%% patterns of effects, e.g., the nesting depth of abstractions in an
%% STLC terms as described in the previous section.  Like ghost
%% variables, ghost operators are used exclusively for specification and
%% verification purposes--= they do not appear in expressions and do not
%% impact the evaluation of a program.

%% \SJ{This example is too detailed.  We should use something simpler -
%%   maybe a stack example that counts pushes and pops?}  \BD{Do we need
%%   to say anything more about ghost events than what we said in the
%%   overview -- our page budget is probably spent on additional examples
%%   of the novel features of the type system.}
%% \begin{example}[Ghost Events]\label{ex:stlc-types}
%%   The \tyName{} specification of the %
%%   $\eff{token}{:}\Code{str} \sarr \Unit$ %
%%   operator from ~\autoref{ex:de-bruijn-index} is:%
%%   \vspace{-.4cm} \par\nobreak {\footnotesize
%%     \begin{align*}
%%       & \effGray{depth}{:}(\Int\sarr\Unit) \garr d{:}\Int \garr
%%       \\&\;\;(t{:}\Code{stlcTy}\garr x{:}\ort{\Code{str}}{\vnu = \boxed{)}} \sarr \uhat{\Code{LAST}(\msgAGray{depth}{d}) \msgA{token}{\boxed{(\lambda t.}\;} \msgAGray{depth}{d + 1} \allA }{\Unit}{\msgA{token}{x} \msgAGray{depth}{d}}) \interty
%%       \\&\;\;(n{:}\Int\garr x{:}\ort{\Code{str}}{\vnu = \boxed{[n]} \land 0 \leq n < d} \sarr \uhat{\Code{LAST}(\msgAGray{depth}{d})}{\Unit}{\msgA{token}{x}})
%%     \end{align*}}\noindent
%%   This \tyName{} consists of two intersected \tyName{} for function
%%   abstraction and bound variable (parameter $x$ is constrained to be
%%   $\boxed{)}$ and $\boxed{[n]}$ respectively), which share the same
%%   ghost operator $\effGray{depth}$ and ghost variable $d$.
%%   \BD{This example needs to be gentler...}
%% \end{example}


% The standard temporal operators
% (e.g., test $\evparenth{\phi}$, next $\nextA A$, until $\untilA$) and
% various set operators (i.e., complement $\neg$, intersection $\land$,
% and union $\lorA$) are defined normally. These operators are
% sufficient to capture other modalities, e.g., eventually ($\finalA$),
% globally ($\globalA$), and importantly, the singleton (last) modality
% $$, which describes a singleton trace, i.e., one which prohibits
% any subsequent effects~\cite{LTLf}. SREs can capture several common
% patterns: the set of all possible traces \(\allTL\), the singleton set
% containing the empty trace \(\epsilon\), and the empty set of traces and
% \(\empTL\); these are analogous to the regular expressions \( .^* \),
% \(\epsilon\), and \(\emptyset\), resp.

% \begin{example}[Queue Library]\label{ex:stack}
% The precise semantics of $\eff{enqueue}$, $\eff{dequeue}$ operation of an effectful queue library requires \emph{first-in-first-out} (FIFO) property which needs counting. We introduce two ghost event $\effGray{numEnq}$ and $\effGray{numDeq}$ which are counters that records the number of previous added and removed elements. Thus, the FIFO property can be defined as ``the $i^{th}$ dequeued element is exactly the $i^{th}$ enqueued element'', and expressed as the following \tyName{}:
% {\footnotesize
%     \begin{align*}
%       \eff{enqueue} :~ & i{:}\Int\garr x{:}\Int\sarr\uhat{\Code{LAST}(\msgAGray{numEnq}{i - 1})}{\msgA{enqueue}{x}\msgAGray{numEnq}{i}}
%       \\\eff{dequeue} :~ & i{:}\Int\garr x{:}\Int\sarr\uhat{\Code{LAST}(\msgA{enqueue}{x}\msgAGray{numEnq}{i}) \land \Code{LAST}(\msgAGray{numDeq}{i - 1})}{\msgA{dequeue}{x}\msgAGray{numDeq}{i}}
%     \end{align*}}\noindent
%   The history regex of $\eff{enqueue}$ operator maintains the counter by get number of previous added elements by match the most last $\effGray{numEnq}$ event in the history (we assume this is number is $0$ when there is no previous $\effGray{numEnq}$) and increase this number by $1$ by appending a new ghost event in future regex. The \tyName of $\eff{dequeue}$ besides maintians the counter $\effGray{numDeq}$ in the same way, also requires the $\eff{dequeue}$ element is exactly the $i^th$ $\eff{enqueue}$ element, i.e., the one update the counter $\effGray{numDeq}$ as $i$.
% \end{example}

\subsection{Typing Rules}

$\DSL{}$ is equipped with a pair of typing judgments on pure and
effectful expressions, $\Gamma \vdash v : t$ and $\Gamma\vdash e:
\uhat{H}{x{:}t}{F}$, respectively. Both judgments rely on a type
context, $\Gamma$, that maps from variables to pure refinement types
(i.e., $t$). As in prior work~\cite{CoverageType}, typing contexts are not
allowed to contain \tyName{}s--- doing so breaks several structural
properties (e.g., weakening) that are used to prove type safety. Both
rules are parameterized over an additional context $\Delta$, which
provides types for pure and effectful operators.

\paragraph{Auxiliary typing relations}
Our type system also relies on three additional auxiliary
relations.\footnote{The full definition of these relations is included
  in the \techreport{}.} The first is a well-formedness relation $\Gamma
\wellfoundedvdash \tau$ that ensures, among other things, that all
qualifiers appearing in a type $\tau$ are closed under the current
typing context $\Gamma$ and that the verification conditions generated by
the synthesis algorithm presented in the next section can be encoded
as effectively propositional (EPR) sentences~\cite{Ramsey1987}, which
can be efficiently handled by an off-the-shelf theorem prover such as
Z3~\cite{de2008z3}.  We also require the
history regex $H$ to not be empty, which guarantees that it represents
a meaningful underapproximate precondition.

\begin{figure}[t!]
  {\footnotesize
  {\footnotesize
  \begin{flalign*}
   &\text{\textbf{Auxiliary Typing}} & & \fbox{$\Gamma \wfvdash \tau \quad \Gamma \vdash A \subseteq A \quad \Gamma \vdash \tau <: \tau$} &\text{\textbf{Typing}} &  & \fbox{$\Gamma\vdash \eff{op}: t \quad \Gamma\vdash v : t \quad \Gamma \vdash e: \tau$}
  \end{flalign*}
  }
  \\ \
  \mprooftr{30mm}
  {$\Gamma\wfvdash \tau[x\mapsto v]$}
  {SubG}
  {$\Gamma \vdash x{:}b\garr \tau <: \tau[x\mapsto v]$}
  \;
  \mprooftr{41mm}
  {$\Gamma \vdash t_1 <: t_2$ \quad
  $\Gamma,x{:}t_1 \vdash F_2 \subseteq F_1$ \quad
  $\Gamma,x{:}t_1 \vdash H_1 \subseteq H_2$
   }
  {SubHF}
  {$\Gamma \vdash \uhat{H_1}{x{:}t_1}{F_1} <: \uhat{H_2}{x{:}t_2}{F_2}$}
  \;
  \mprooftr{30mm}
  {$\Gamma \vdash \eff{op} : \overline{x{:}t_x}\sarr\uhat{H}{t}{F}$ \quad
  $\Gamma,\overline{x{:}t_x} \vdash H' \subseteq H$ \quad
  $\Gamma,\overline{x{:}t_x} \not\vdash H' \subseteq \emptyset$}
  {TOpHis}
  {$\Gamma\vdash \eff{op} : \overline{x{:}t_x}\sarr\uhat{H'}{t}{F}$}
  % \quad
  % \mprooftr{30mm}
  % {$\Gamma\vdash e: \tau$ \quad
  % $\Gamma \vdash \tau <: \tau'$\quad $\Gamma \vdash \tau' <: \tau$}
  % {TEq}
  % {$\Gamma\vdash e: \tau'$}
  \\[.8em]
  \mprooftr{20mm}
  {$\Gamma\vdash v : t$}
  {TRet}
  {$\Gamma\vdash v : \uhat{H}{t}{\epsilon}$}
  \quad
  \mprooftr{15mm}
  {$\Delta(\eff{op}) = t$}
  {TOpCtx}
  {$\Gamma\vdash \eff{op} : t$}
  \quad
  \mprooftr{58mm}
  {
  $\Gamma \vdash \eff{op} : \overline{x_i{:}t_i}\sarr \uhat{H}{t}{\msgA{op}{\overline{x_i}} \seqA F}$ \quad
  $\forall i. \;\Gamma \vdash v_i : t_i$}
  {TEffOp}
  {$\Gamma\vdash \eff{op}{\;\overline{v_i}} : \uhat{H}{t}{\msgA{op}{\overline{v_i}} \seqA F}$}
  \\[.8em]
  %%\mprooftr{56mm}
  %%{
  %%$\Gamma \vdash \eff{op} : \overline{x_i{:}t_i}\sarr \uhat{H}{t}{\msgA{op}{\overline{x_i}} \seqA F}$ \\
  %%  $\forall i. \;\Gamma \vdash v_i : t_i$}
  %% {TReqOp}
  %% {$\Gamma\vdash \reqOp{op}{\;\ensuremath{\iota\; \overline{v_i}}} : \uhat{H}{\Code{id}}{\reqEvt{op}{\ensuremath{\iota\; \overline{v_i}}} \seqA F}$}
  %% % \mprooftr{62mm}
  %% % {
  %% % $\Gamma\vdash \eff{op} : \overline{x_i{:}t_i}\sarr \uhat{H}{\Unit}{\msgA{op}{\overline{x_i}} \seqA F}$ \quad
  %% % $\forall i. \;\Gamma \vdash v_i : \widehat{t_i}$ \quad
  %% % $\Gamma\vdash e : \uhat{H \seqA (\msgA{op}{\overline{v_i}} \land \msg{op}{\phi}) \seqA F}{t}{F'}$}
  %% % {TGen}
  %% % {$\Gamma\vdash \zgen{op}{\overline{v_i}}{e} : \uhat{H}{t}{\msg{op}{\phi} \seqA F \seqA F'}$}
  %% \;
  %% \mprooftr{51mm}
  %% {
  %%   $\Gamma \vdash \eff{op} : \overline{x_i{:}t_i}\sarr \uhat{H}{t}{\msgA{op}{\overline{x_i}} \seqA F}$\\
  %%   $\forall i. \;\Gamma \vdash v_i : t_i$
  %% }
  %% {TRspOp}
  %% {$\Gamma\vdash \respOp{op}{\ensuremath{\iota}} : \uhat{H}{t}{\respEvt{op}{\ensuremath{\iota\; \overline{v_i}}} \seqA F}$}
  %% \\[.8em]
  \mprooftr{35mm}
  {
  $\Gamma\vdash e_1 : \tau_1$ \quad $\Gamma\vdash e_2 : \tau_2$}
  {TChoice}
  {$\Gamma\vdash e_1 \oplus e_2 : \tau_1 \interty \tau_2$}
  \quad
  \mprooftr{66mm}{
    $\Gamma\vdash e_1 : \uhat{H}{x{:}t_x}{F_1}$ \quad
    $\Gamma, x{:}t_x\vdash e_2 : \uhat{H \seqA F_1}{t}{F_2}$}
  {TLet}
  {
    $\Gamma\vdash \zlet{x}{e_1}{e_2} : \uhat{H}{t}{F_1 \seqA F_2}$
  }
  \\[.8em]
  \mprooftr{22mm}{
   $\Gamma, x{:}t_x\vdash e : \tau$
  }{TFun}{
   $\Gamma\vdash \zzlam{x}{e} : x{:}t_x\sarr \tau$}
  \
  \mprooftr{47mm}{
  }{TFixBase}{
   $\Gamma\vdash \zzfix{f}{x}{e} : \overline{x{:}s}\sarr\uhat{H}{\urt{b}{\bot}}{\emptyset}$}
  \
  \mprooftr{27mm}{
   $\Gamma\vdash \zzfix{f}{x}{e} : t$ \quad
   $\Gamma, f{:}t\vdash \zzlam{x}{e} : t'$
  }{TFixInd}{
   $\Gamma\vdash \zzfix{f}{x}{e} : t'$}
  % \quad
  % \mprooftr{60mm}{
  %  $\Gamma\vdash v_1 : y{:}t_y\sarr\uhat{H}{A}$ \quad
  %  $\Gamma\vdash v_2 : t_y$ \quad
  %  $\Gamma; \Delta' \vdash e' : \uhat{H \seqA A}{A'}$
  % }{TApp}{
  %  $\Gamma\vdash \zlet{x}{v_1\;v_2}{e} : \uhat{H}{A \seqA A'}$}
  }
  \vspace*{-.1in}
  \caption{Selected typing rules.}
  \label{fig:selected-typing-rules}
  \vspace*{-.15in}
  \end{figure}

  The type system also uses a semantic subtyping relation tailored for
  underapproximation; \autoref{fig:selected-typing-rules} presents the
  key subtyping rule for \tyName{}s.  The rule \textsc{SubG} allows
  the instantiation of a ghost variable in a \tyName{} into
  well-formed types. The rule \textsc{SubHF} uses an auxiliary
  inclusion relation on SREs $\Gamma \vdash A \subseteq A$ to relate
  the history and future regexes of two \tyName{}s under the current
  type context $\Gamma$, as well as the return type using the standard
  subtyping relation on pure refinement types. Notably, because our
  typing judgements are underapproximate, this subtyping relation is
  covariant in the history regex and contravariant in the future
  regex, yielding behavior similar to the reverse rule of consequence
  found in Incorrectness Logic (IL)~\cite{OHP19}.  Like that rule,
  \textsc{SubHF} allows us to strengthen a postcondition (i.e., shrink
  the future trace), and weaken the precondition (i.e., enlarge the
  history trace).  Since the future automata in a \tyName{} can refer
  to the pure value it returns, the inclusion check between the future
  automata extends the typing context with a binding for the return
  value.

A subset of our typing rules is shown in
\autoref{fig:selected-typing-rules}.\footnote{The complete set of
  typing rules is included in the supplemental material.}  All our
rules assume any types they use are well-formed, so we elide the
corresponding well-formedness judgments from their premises.  The
\textsc{TRet} rule requires the future regex of a pure term
($\epsilon$) to only accept the empty trace.  The operator type is
retrieved from the global operator context $\Delta$ using the
\textsc{TOpCtx} rule.  For a given event, there may be multiple
reachable history traces (underapproximate preconditions) that can be
allowed by an operator that is determined by the context in which the
operator executes, and by the operator type alone.  We, therefore,
expect that the operator type found in the operator context is the
most precise one, i.e., it doesn't allow any unreachable history
traces, and admits underapproximation specialization.  The rule
\textsc{TOpHis} can then freely choose a \emph{non-empty} subset of the reachable
history traces to align with the actual calling context.  The
application rule \textsc{TEffOp} requires the arguments of an
effectful operator to be well-typed against its parameter types; it
records the effect as an event in the future trace along with any
additional ghost events ($F$).  The nondeterministic choice operator
is typed using the \textsc{TChoice} rule which merges the \tyName{}s
of the two branches.  The \textsc{TLet} rule requires the binding
expression $e_1$ to generate the first part of its future trace $F_1$,
with the rest of trace produced by the body expression $e_2$ under a
new history that includes $F_1$.  The typing rule for functions
(\textsc{TFun}) is standard, while the rule for recursive functions is
modeled after the corresponding IL rule, unrolling the recursion as
needed.  The \textsc{FixBase} rule types any recursive function with a
\tyName{} lacking reachability constraints (i.e.,
$\uhat{H}{\urt{b}{\bot}}{\emptyset}$). In contrast, the
\textsc{FixInd} rule permits typing a recursive function with type
$t'$ by requiring it to have another type $t$, and by checking that
its body can be assigned type $t'$ under the assumption that the
function itself has type $t$.

%% The \textsc{TReqOp} and the
%% \textsc{TRspOp} rule handle asynchronous effects, using
%% well-formedness conditions to ensure that the values carried by a
%% response are the same as those of a request in the history trace with
%% the same id.

\begin{example}[Operator Application Typing]
  \label{ex:operator-application-typing}
  We present a typing derivation for a simple generator that generates
  $\eff{push}$ and $\eff{pop}$ operations using the signatures from
  Example~\ref{ex:stack-example}. Consider how we might type check the
  following expression: %
  \vspace{-.4cm} \par\nobreak {\small %
    \begin{align*}
      &\emptyset \vdash \zlet{x{:}\Code{int}}{\sassume{0 < x}}{\eff{push}\;x; \eff{pop}} :
      \\&\qquad \uhat{\msgAGray{pushI}{0\;0}\msgAGray{popI}{0\;0}}{x{:}\urt{\Int}{0 < \vnu}}{\msgA{push}{x}\msgAGray{pushI}{1\;x}\msgA{pop}{x}\msgAGray{popI}{1\;x}}
    \end{align*}}%
  \noindent
  The signature of the generator is a \tyName{} that assumes a history
  trace in which the stack is empty (i.e., the ghost events
  $\effGray{pushI}$ and $\effGray{popI}$ both hold a $0$ index and a
  dummy value).  The body of the generator generates a positive number
  $x$ using $\mykw{assume}$; the underapproximate refinement type of
  $x$ signifies that it must evaluate to every positive integer. The
  generator then is expected to pop the same value.  We can thus type
  check its two effectful operation applications using the following
  judgement: %
  \vspace{-.4cm} \par\nobreak {\small %
    \begin{flalign*}
      &x{:}\urt{\Int}{0 < \vnu} \vdash \eff{push}\;x; \eff{pop} :
      \\&\qquad\qquad\ \uhat{\msgAGray{pushI}{0\;0}\msgAGray{popI}{0\;0}}{x{:}\urt{\Int}{0 < \vnu}}{\msgA{push}{x}\msgAGray{pushI}{1\;x}\msgA{pop}{x}\msgAGray{popI}{1\;x}}
    \end{flalign*}}%
  \noindent
  Using the \textsc{TEffOp} rule to type the first $\eff{push}$
  effect, we retrieve the operator's types from the operator context
  $\Delta$:%
  \vspace{-.4cm} \par\nobreak {\small %
    \begin{flalign*}
      &... \vdash \eff{push} : n{:}\Int \garr y{:}\Int \garr &&
      \\[-0.2em]& \qquad\qquad\ \  x{:}\urt{\Int}{y < \vnu} \sarr
      \uhat{\Code{Last}(\msgAGray{pushI}{n\;y})}{\Unit}{\msgA{push}{x}\msgAGray{pushI}{n{+}1\;x}}  &&\text{(by \textsc{TOpCtx})}
    \end{flalign*}}\noindent We can then apply the subtyping rule
  (\textsc{SubG}) to instantiate the signature's ghost variables
  (e.g., $n \mapsto 0$ and $y \mapsto 0$): \vspace{-.4cm} \par\nobreak
  {\small %
    \begin{flalign*}
      &... \vdash \eff{push} : x{:}\urt{\Int}{0 < \vnu} \sarr
      \uhat{\Code{Last}(\msgAGray{pushI}{0\;0})}{\Unit}{\msgA{push}{x}\msgAGray{pushI}{1\;x}} &&\text{(by \textsc{SubG})}
    \end{flalign*}}%
  \noindent
  The rule \textsc{TOpHis} can then be used to specialize its history
  trace to align with that of the actual calling context
  ($\msgAGray{pushI}{0\;0}\msgAGray{popI}{0\;0}$). Note that we do not
  need to also specialize the future regex, since it is already a
  prefix of the future trace of the target type: %
  \vspace{-.4cm} \par\nobreak {\small %
    \begin{flalign*}
      &... \vdash \eff{push} : x{:}\urt{\Int}{0 < \vnu} \sarr
      \uhat{\msgAGray{pushI}{0\;0}\msgAGray{popI}{0\;0}}{\Unit}{\msgA{push}{x}\msgAGray{pushI}{1\;x}}
      &&\text{(by \textsc{TOpHis})}
    \end{flalign*}}%
  \noindent
  The rest of program is now typed via the following judgement:
  \vspace{-.4cm} \par\nobreak {\small %
    \begin{align*}
      &x{:}\urt{\Int}{0 < \vnu}\vdash \eff{pop} : \uhat{\msgAGray{pushI}{0\;0}\msgAGray{popI}{0\;0}\msgA{push}{x}\msgAGray{pushI}{1\;x}}{x{:}\urt{\Int}{0 < \vnu}}{\msgA{pop}{x}\msgAGray{popI}{1\;x}} &&
    \end{align*}}%
  \noindent
  As before, after looking up the type of $\eff{pop}$ using
  \textsc{TOpCtx} (\ref{eqn:typing+l1}), we can instantiate the ghost variable $n$ as $1$
  and $m$ as $0$ via \textsc{SubG} (\ref{eqn:typing+l2}), align the return type via
  \textsc{SubHF} ((\ref{eqn:typing+l3}) and refine the history regex via the rule
  \textsc{TOpHis} (\ref{eqn:typing+l4}): %
  \vspace{-.4cm} \par\nobreak {\footnotesize %
    \begin{flalign}
      &... \vdash \eff{pop} : n{:}\Int \garr m{:}\Int \garr \notag
      \\[-0.2em]& \qquad\qquad\  \uhat{\Code{Last}(\msgAGray{popI}{m\;\_}) \land
      \Code{Last}(\msgAGray{pushI}{n\;\_}) \land
      (\allA\msgAGray{pushI}{(n{-}m)\;x}\allA)}{x{:}\Int}{\msgA{pop}{x}\msgAGray{popI}{m{+}1\;x}}
    && \tag*{\ding{172}}\label{eqn:typing+l1}
      \\&... \vdash \eff{pop} : \uhat{\Code{Last}(\msgAGray{popI}{0\;\_}) \land
      \Code{Last}(\msgAGray{pushI}{1\;\_}) \land
      (\allA\msgAGray{pushI}{1\;x}\allA)}{x{:}\Int}{\msgA{pop}{x}\msgAGray{popI}{1\;x}}
    && \tag*{\ding{173}}\label{eqn:typing+l2}
      \\&... \vdash \eff{pop} : \uhat{\Code{Last}(\msgAGray{popI}{0\;\_}) \land
      \Code{Last}(\msgAGray{pushI}{1\;\_}) \land
      (\allA\msgAGray{pushI}{1\;x}\allA)}{x{:}\urt{\Int}{0 <
        \vnu}}{\msgA{pop}{x}\msgAGray{popI}{1\;x}} &&
    \tag*{\ding{174}}\label{eqn:typing+l3}
      \\&... \vdash \eff{pop} :
      \uhat{\msgAGray{pushI}{0\;0}\msgAGray{popI}{0\;0}\msgA{push}{x}\msgAGray{pushI}{1\;x}}{x{:}\urt{\Int}{0
          < \vnu}}{\msgA{pop}{x}\msgAGray{popI}{1\;x}} &&
      \tag*{\ding{175}}\label{eqn:typing+l4}
    \end{flalign}}%
  \noindent
\end{example}

% \begin{example}[Operator Type Context] The type of effect operator
%   $\eff{readRsp}$ in Example~\ref{ex:read-atomicity} can provided by
%   the following operator context $\Delta$: %
%   \vspace{-.4cm} \par\nobreak {\small %
%     \begin{align*}
%       \Delta \equiv \{&(\eff{readRsp}, \iota{:}\Int\sarr \uhat{\msgA{writeRsp}{\iota\;x}\msgA{readReq}{\iota}}{x{:}\urt{\Int}{\top}}{\msgA{readRsp}{\iota\;x}}),\\
%                       &(\eff{readRsp}, \iota{:}\Int\sarr
%                         \uhat{\msgA{writeRsp}{\iota\;x}\msgA{readReq}{\iota}\msgC{writeRsp}}{x{:}\urt{\Int}{\top}}{\msgA{readRsp}{\iota\;x}}),
%                         ... \}
%     \end{align*}}
% \end{example}

% The typing rules effectful operations, \textsc{TGen} and
% \textsc{TObs}, both extract the operator type
% $\uhat{H}{t}{\msg{op}{\overline{x}}\seqA F}$ from operator context
% $\Delta$ \BD{The current rules do not mention $\Delta$, need to fix};
% and require that it aligns with the \tyName{} of the expression that
% the operation is being performed in:%
% The rules for performing events, \textsc{TGen} and \textsc{TObs}, both extract
% the type of the corresponding operator from $\Delta$,
% $\rg{H}{\msg{op}{\phi}}{A \seqA F}$, and require that it aligns
% with the \tyName{} of the expression that the operation is being
% performed in:%
% \vspace{-.4cm} \par\nobreak {\footnotesize %
%   \begin{align*}
%     \text{\textcolor{gray}{(term's \tyName)}} \qquad
%     \effLR{\underbrace{ H}_\texttt{history}}\;t\;
%     \effLR{
%     \;\;\underbrace{\msg{op}{\overline{x}}}_\texttt{trace produced by $\eff{op}$}
%     \;\;\underbrace{F }_\texttt{ghost events appended to $\eff{op}$}
%     \;\;\underbrace{F'}_\texttt{trace produced by $e$}
%     }
% \end{align*}}%
% \SJ{I don't follow what it means to append a ghost event to a normal
%   event.}
% \noindent For simplicity, we only allow the ghost events to be
% appended to a normal event.  Same as prior work~\cite{CoverageType},
% the arguments should be well-typed against the parameter types fliped
% from over- to under-approximation (i.e., $\widehat{t_i}$).  To type
% the rest of the expression, the assumption of \tyName of operator
% (i.e., $\msg{op}{\overline{x}}$) should be aligned with the
% expectation of term's type (i.e., $\msg{op}{\phi}$), then both rules
% move the symbolic event
% $(\msg{op}{\overline{x}} \land \msg{op}{\phi})\seqA F$ from the future
% regex to the tail of the history regex. The subsumption rule
% \textsc{TSub} and the equivalence rule \textsc{TEq} allows us to
% change the shape of a type that a term is being typed against.


% The nondeterministic loop aligns with the Kleene star in SRE, and thus can be straightforwardly typed by the \textsc{TLoop} rule.
% Since the loop body $e$ may be executed an unbounded number of times, the typing judgment for $e$ requires that an unbounded number of iterations may have already been performed ($H \seqA \textcolor{orange}{A^*}$) and may occur in the future ($\textcolor{orange}{A^*} \seqA A' \seqA F$). Additionally, the input and output capability contexts must be identical, ensuring that $e$ replenishes any capabilities it consumes.

% of observing effect operator $\eff{op}$ from the typing
% context with the help of effect operator typing
% $\textsc{TEOp}$. Besides requiring the arguments ($\overline{v_i}$) is
% well-typed against parameter types ($\overline{t_i}$) as application
% rules in standard refinement typing, it additionally

\subsection{Metatheory}
\label{sec:metatheory}


% \paragraph{Operational Semantics} Events in $\DSL$ are operations
% applied to concrete values $(\zevent{op}{\overline{c}})$. Evaluating a
% $\DSL$ program depends on an context \emph{trace}, i.e., a sequence of
% events, and an \emph{observation capability}, i.e., a multiset of
% generable events that are waiting to be handled by observable
% operators. Each evaluation step produces an output trace and an
% updated capability. Traces are equipped with the standard list
% operations (i.e., cons ${::}$ and concatenation $\listconcat{}$).  The
% operational semantics of $\DSL$ are defined by the small-step
% reduction relation:
% $\steptr{\alpha}{(\beta, e)}{~\alpha'}{(\beta', e')}$.  This judgment
% is read as: ``under the history trace $\alpha$ and observation
% capability $\beta$, $e$ steps to $e'$, emitting the trace $\alpha'$
% and remaining the capability $\beta'$.'' Intuitively, the history
% trace $\alpha$ represents the sequence of events visible to a
% operator, thereby determining its response; the capability $\beta$
% contains events that are able to be observed in the future. The
% semantics uses an auxiliary judgement,
% $\alpha \vDash m \Downarrow m'$, that specifies the observable event
% $m'$ that will be generated after handling a generative event $m$.

% \autoref{fig:selected-semantics} presents the key rules of $\DSL$'s
% semantics,\footnote{The remaining rules are completely standard and
%   provided in the supplemental material.} which reflect the
% ``consume-and-produce'' semantics of observation capability.  The rule
% for observable events (\textsc{StObs}) removes an event $m$ from the
% capability that can produce the observable operation $\eff{op}$,
% evaluates it under the current context, and substitutes the observed
% value $c'$ (last payload of $m$) for binding variable $x$ in $e$, the
% body of the $\mykw{let}$ expression.  In contrast, the rule for
% generable events (\textsc{StGen}) makes no assumptions on the
% observation capability, since these events can be directly performed
% by the generator. It simply add itself to the resulting observation
% capability for future handling.

% \begin{example}[Auxiliary Operator Semantics]\label{ex:aux-semantics}
%   The auxiliary semantics of the $\eff{read}$ and $\eff{write}$ operators in Example~\ref{ex:read-atomicity} can be defined as follows:

%   {\footnotesize
%     \mprooftr{44mm}{}{OpWrite}{
%     $\alpha \vDash \zevent{writeReq}{\iota,v} \Downarrow \zevent{writeRsp}{\iota,v}$
%     }
%     \quad
%     \mprooftr{62mm}{$\alpha_1 = \alpha_{11} \listconcat \zevent{writeRsp}{\iota',v} \listconcat \alpha_{12}$ \quad $\eff{readRsp}\not\in \alpha_{12}$}{OpRead}{
%     $\alpha_1\listconcat\zevent{readReq}{\iota}\listconcat\alpha_2\vDash \zevent{readReq}{\iota} \Downarrow \zevent{readRsp}{\iota,v}$
%     }
%     \\[-0.1em]
%     }\noindent

%   \noindent
%   The database responds to a $\eff{writeReq}$ event by emitting a
%   $\eff{writeRsp}$ event containing the identical payload. To ensure
%   read atomicity, handling a $\eff{readReq}$ event requires
%   identifying the corresponding history prefix $\alpha_1$ at the point
%   of the read request happens, and producing a $\eff{readRsp}$ event
%   whose value $v$ matches the most recently written value for that key
%   in $\alpha_1$.
% \end{example}

% \begin{example}[Read Atomicity]\label{ex:semantics-async}
%   The following derivation illustrates the execution of a single iteration of the program from Example~\ref{ex:read-atomicity}, evaluated under the auxiliary semantics for $\eff{writeReq}$ and $\eff{readReq}$ specified in Example~\ref{ex:aux-semantics}, which enforce read atomicity. For this iteration, we assign transaction identifiers $\Code{tid1}=1$ and $\Code{tid2}=2$, and two random generated integers are $7$ and $8$:
%     {\footnotesize
%       \begin{flalign*}
%         &[] \vDash
%         (\emptyset, \eff{writeReq}\;1\;7;\zobs{\_}{writeRsp}{1}{\eff{readReq}\;2; \textit{line 5 - 6 of code in Example~\ref{ex:read-atomicity}}}) &&
%         \\[-0.2em]
%         &\;\;\xhookrightarrow{[\zevent{writeReq}{1,7}]}{(\{ \zevent{writeReq}{1,7} \}, \zobs{\_}{writeRsp}{1}{\eff{readReq}\;2; \textit{line 5 - 6 of code in Example~\ref{ex:read-atomicity}}})}  \tag{\textsc{StGen}} &&
%         \\[-0.4em]
%         &\;\;\xhookrightarrow{[\textcolor{blue}{\zevent{writeRsp}{1,8}}]}{(\emptyset, \eff{readReq}\;2; \eff{writeReq}\;1\;7;\zobs{\_}{writeRsp}{1}{\textit{line 6 of code in Example~\ref{ex:read-atomicity}}})}  \tag{\textsc{StObs}} &&
%         \\[-0.4em]
%         &\;\;\xhookrightarrow{[\textcolor{orange}{\zevent{readReq}{2}]}}{(\{\zevent{readReq}{2}\}, \eff{writeReq}\;1\;8;\zobs{\_}{writeRsp}{1}{\textit{line 6 of code in Example~\ref{ex:read-atomicity}}})}  \tag{\textsc{StGen}} &&
%         \\[-0.4em]
%         &\;\;\xhookrightarrow{[\zevent{writeReq}{1,8}]}{(\{\zevent{readReq}{2}, \zevent{writeRsp}{1,8}\}, \zobs{\_}{writeRsp}{1}{\textit{line 6 of code in Example~\ref{ex:read-atomicity}}})}  \tag{\textsc{StGen}} &&
%         \\[-0.4em]
%         &\;\;\xhookrightarrow{[\zevent{writeRsp}{1,8}]}{(\{\zevent{readReq}{2}\}, \zobs{y}{readRsp}{2}{\sassert{y == 7}; () \oplus \Code{transaction\_gen\;1\;2}})}  \tag{\textsc{StObs}} &&
%         \\[-0.4em]
%         &\;\;\xhookrightarrow{[\zevent{readRsp}{2,7}]}{(\emptyset, \sassert{7 == 7}; () \oplus \Code{transaction\_gen\;1\;2})}  \tag{\textsc{StObs}} &&
%       \end{flalign*}
%     }\noindent
%   In the final step, the program executes $\mykw{obs}\;\eff{readRsp}\;2$, producing a read value $y=7$, which matches the value of the most recent $\eff{writeRsp}$ (blue) at the point where the $\eff{readReq}$ event (orange) occurred. Thus, read atomicity is preserved and execution can continue to generate further transactions.
% \end{example}

Although ghost events constrain \tyName{}s, they do not
appear in actual traces produced by the operational semantics. We
define an erasure function $\eraseGhost{\alpha}$ to remove ghost events from a trace $\alpha$.

\paragraph{Type denotations} Similar to other refinement type
systems~\cite{JV21}, types in $\DSL$ are denoted as their inhabitants
(i.e., $\denot{t}$ and $\denot{\tau}$).  The type context $\Gamma$ is
denoted as a \emph{substitution} $\sigma$ (i.e.,
$[x_1\mapsto v_1, x_2\mapsto v_2]$) that provides the assignments for
binding variables in $\Gamma$ and the denotation (the denoted
language) of SREs is a set of traces. Then, trace inclusion under a
type context is defined as
$\Gamma \vdash A \subseteq A' \doteq \forall \sigma \in
\denot{\Gamma}. \denot{\sigma(A)} \subseteq \denot{\sigma(B)}$.  The
type denotation of $\uhat{H}{x{:}\urt{b}{\phi}}{F}$ is:
\vspace{-.4cm} \par\nobreak {\small\begin{align*}
  \{e ~|~ \emptyset \basicvdash e : b \land \forall v{:}b. \phi[v\mapsto v] \impl \forall \alpha_f \in \denot{F[x\mapsto v]}. \exists \alpha_h \in\denot{H[x\mapsto v]}. \msteptr{\eraseGhost{\alpha_h}}{e}{\eraseGhost{\alpha_f}}{v} \}
\end{align*}}\noindent
A term $e$ inhabits $\uhat{H}{x{:}\urt{b}{\phi}}{F}$ iff, for trace $\alpha_f$ accepted by $F$, there exists history trace $\alpha_h$ accepted by $H$, such that the execution of $e$ under
$\alpha_h$ produces $\alpha_f$.
\footnote{The details of these definition can be found in our \techreport{}.}
Our formulation constitutes a trace-based analogue of IL~\cite{OHP19}, in which effects are modeled as program state transitions:
\vspace{-.4cm} \par\nobreak {\small %
  \begin{align*}
    [P]\; e\; [Q] \doteq \forall s_\Code{post} \in \denot{Q}. \exists s_\Code{pre} \in \denot{P}. (s_\Code{pre}, e) \Downarrow s_\Code{post}
  \end{align*}}\noindent
A notable distinction between these two presentations is that we
characterize $F$ as the set of \emph{newly produced} traces rather
than the full post-execution trace, reflecting our trace-based
type rules. This formulation additionally avoids
redundant prefixes inherited from the execution history.

% In contrast to incorrectness logic~\cite{OHP19}, which explicitly models effects via program states, our latent effect semantics require executions to terminate with an empty capability context. Permitting termination with a non-empty capability (e.g., $\{\eff{readRsp}\}$) would leave latent effects from asynchronous actions (such as $\eff{readReq}$) unobserved. Thus, we explicitly denote the \tyName with term and its continuation $e_f$ to capture the entire execution. The $\beta$ and $\beta'$ are the capability before and after execution of $e$; since it is always valid to add unused capability to both sides, thus this denotation is parameterized by capability context $\Theta_A$ and $\Theta_F$.

% Underapproximating specifications, as in our approach, must exclude unreachable states or unrealizable traces, since $A$ cannot include them. To improve expressiveness and usability, our tool allows users to specify an overapproximated style global property $P$ and verifies that $A \subseteq P$.

%  A term $e$ inhabits $\rg{H}{A}{F}$ iff, for all realizable (i.e., producible by the execution of some program) history traces $\alpha_h$ accept by $H$, the current trace $\alpha$ produced by $e$ is accepted by $A$, and all future traces $\alpha_f$ accepted by $F$ are also realizable.\footnote{The details of these definition can be found in our \techreport{}.}

% \paragraph{Operator context}
% We expect the trace produced by our test generator to be consistent with the operator context $\Delta$; i.e., each event in the trace must be consistent with corresponding $\tyName$ in $\Delta$.

% \begin{definition}[Trace consistent with operator context]
%   \label{lemma:trace-conistency}
%   A trace $\alpha$ is consistent with a operator context $\Delta$
%   (written $\traceDelta{\Delta}{\alpha}$) iff for every
%   $\zevent{op}{\overline{c_i}}$ in $\alpha$ (i.e., $\alpha = \alpha_h
%   \listconcat [\zevent{op}{\overline{c_i}}] \listconcat \alpha_f$),
%   there exists $\overline{x{:}t} \sarr \rg{H}{\msg{op}{\phi}}{F}$ and
%   capability $\Theta$ in $\Delta$ such that the following hold: %
%   \vspace{-.4cm} \par\nobreak {\small %
%     \begin{align*}
%     \alpha_h \in \denot{H\msubst{x_i}{c_i}} \land \alpha_f \in \denot{F\msubst{x_i}{c_i}} \land \Theta \subseteq \{ \eff{op} ~|~ \zevent{op}{\overline{c}} \in \alpha_f\} \land \phi\msubst{x_i}{c_i} \land \forall i. c_i \in \denotation{t_i}
% \end{align*}}
% \end{definition}

We expect the \tyName{}s provided by the operator typing context
$\Delta$ to be consistent with the corresponding auxiliary operator
semantics:
\begin{definition}[Well-formed operator typing context]
  \label{lemma:built-in-typing}
  The operator typing context $\Delta$ is well-formed
  iff the semantics of every operator $\eff{op}$ is
  consistent with the type $\overline{x{:}b_{\I{x}}}\garr
  \overline{y{:}t_{\I{y}}}\sarr \uhat{H}{z{:}t}{F}$ provided by $\Delta$, and $H$
  does not specify any unreachable traces.
  \vspace{-.4cm} \par\nobreak {\small\begin{align*}
\overline{\forall
      x{:}b_{\I{x}}}.\; \overline{\forall y \in \denotation{t_{\I{y}}}}.\; \forall H'. H'\not=\emptyset \land  H' \subseteq H \impl (\eff{op}\, \overline{y}) \in \denotation{\uhat{H'}{z{:}t}{F}}
  \end{align*}} \vspace{-.4cm} \par\nobreak\noindent
\end{definition}

\begin{theorem}[Fundamental Theorem]
  \label{theorem:fundamental} A well-typed term, i.e.,
  \(\Gamma\vdash e : \uhat{H}{t}{F}\), generates traces consistent
  with its \tyName:
  $\forall \sigma \in \denotation{\Gamma}. \sigma(e) \in
  \denotation{\sigma(\uhat{H}{t}{F})}$.\footnote{The proofs of all
    theorems in this paper are provided in the \techreport{}.}
\end{theorem}

\begin{corollary}[Type Soundness]
  \label{theorem:sound} A generator \(e\) that satisfies
  \(\emptyset \vdash e : \uhat{\epsilon}{\Unit}{F}\) must produce all
  traces accepted by $F$, i.e.,
  $\forall \alpha \in
  \denot{F}. \msteptr{[]}{e}{\eraseGhost{\alpha}}{()}$.
\end{corollary}
